\section{题08：魔法前缀（P16790）}

\subsection{问题重述}

给定 $N$ 个小写字母单词，定义有效前缀集合为所有单词的全部非空前缀的集合（相同前缀去重）。小蓝和小桥轮流操作，小蓝先手。每次操作时，当前玩家从集合中选择一个非空前缀 $P$，然后删除集合中所有以 $P$ 为前缀的字符串（包括 $P$ 本身）。轮到某玩家时若集合为空则该玩家判负。两人均采取最优策略。每组测试用例判断先手小蓝（XiaoLan）还是后手小桥（XiaoQiao）必胜。约束：$T \le 10$，$N \le 10^5$，所有测试用例字符串总长度不超过 $10^6$。

\subsection{模型转化：Trie树上的公正组合游戏}

第一步是将字符串集合转化为Trie树。将所有单词插入一个Trie树，根结点（编号0）代表空前缀，不出现在游戏中；每个非根结点唯一对应一个有效前缀。于是游戏操作变为：选择任意一个非根结点，删除该结点及其整棵子树。被删除的子树从游戏中移除，但剩余部分可能分裂成多个互不影响的独立子树，每个子树以原根结点的某个未被删除的子结点为根。

例如，单词 \texttt{cat} 和 \texttt{car} 的Trie树包含：根结点（空），根的子结点 \texttt{c}，\texttt{c} 的子结点 \texttt{a}，\texttt{a} 的子结点 \texttt{t} 和 \texttt{r}。若某玩家选择结点 \texttt{a}，则 \texttt{a}、\texttt{t}、\texttt{r} 均被删除，仅剩 \texttt{c} 挂在根下。若该玩家选择的是 \texttt{c}，则整棵树被清空，游戏即刻结束。

这是一个公正组合游戏：双方可进行的操作完全相同，信息完全公开，胜负仅取决于状态而非玩家身份。根据Sprague-Grundy定理，每个状态对应一个Grundy数，终态（无任何结点的空集）Grundy为0；非终态的Grundy数等于其所有后继状态Grundy数的mex（最小非负未出现整数）。整个游戏由根结点的各个子结点为根的子树作为独立子游戏构成，总Grundy数等于各子游戏Grundy数的异或（Nim加法）和。若总Grundy数非零则先手必胜，否则后手必胜。

\subsection{朴素思路：穷举博弈树为何不可行}

最直接的想法是将每个有效前缀集合视为一个博弈状态，枚举当前玩家可选择的所有前缀，模拟双方最优走法。但状态数量随前缀集合大小呈指数增长：假设Trie有 $M$ 个非根结点，每次操作删除一个子树后剩余部分仍可能有多种组合方式，状态空间远远大于 $2^M$。而本题中总结点数可达 $10^6$ 级别，穷举博弈树完全不可行。

\subsection{子树Grundy递归公式的核心洞察}

关键突破在于观察操作的局部性质。在Trie树上，选择某个结点 $v$ 仅影响 $v$ 的子树内部结构，$v$ 的兄弟子树完全不受影响。换言之，对 $v$ 子树的操作是一个独立的子游戏，其Grundy数仅取决于 $v$ 下方树的结构，与整棵Trie的其余部分无关。

考虑以某个非根结点 $v$ 为根的子树。设 $v$ 有 $k$ 个子结点 $c_1, c_2, \dots, c_k$，以它们为根的子树的Grundy数分别为 $g_1, g_2, \dots, g_k$。我们想要求出 $v$ 的子树整体对应的Grundy数 $G(v)$。

当前玩家面对 $v$ 这棵子树时有两种选择：一是直接选择 $v$ 本身，将整棵子树清空，后继状态Grundy为 $0$。二是在某个子结点 $c_i$ 的子树内部进行操作，子游戏 $c_i$ 的Grundy数从 $g_i$ 变为其某个后继值 $h$（由Grundy数的定义，$c_i$ 的所有后继Grundy数恰好覆盖 $0, 1, \dots, g_i - 1$ 中至少这些值），而其他子游戏保持不变。此时整体Grundy变为 $1 + h \oplus \bigoplus_{j \ne i} g_j$。这里加 $1$ 是因为对 $c_i$ 子树的操作发生于 $v$ 的子树内部，$v$ 本身仍存在于游戏中；从博弈代数角度看，$v$ 与 $k$ 个子游戏之间是“串联”关系：$v$ 的一次额外选择权叠加在 $k$ 个并列子游戏之上。

可以严格证明 $G(v) = 1 + \bigoplus_{i=1}^k g_i$。下面用一个三结点小例来展示直观含义：考虑一条长度为 $2$ 的链 \texttt{root} -- \texttt{a} -- \texttt{b}。叶结点 \texttt{b} 无子结点，$G(\texttt{b}) = 1 + 0 = 1$。结点 \texttt{a} 有一个子结点 \texttt{b}（其 $g = 1$），故 $G(\texttt{a}) = 1 + 1 = 2$。这意味着：若玩家直接选 \texttt{a}，游戏结束（Grundy 0）；若玩家选 \texttt{b}，则 \texttt{a} 仍在场，剩下一个Grundy为 $1$ 的子树，后继Grundy为 $1$。可达集合为 $\{0, 1\}$，mex 正好是 $2$。公式成立。

\subsection{算法推导：自底向上计算总Grundy值}

算法流程如下。首先建Trie：将所有单词插入一棵Trie树，每个结点维护26个子结点指针。由于Trie的建树顺序保证子结点编号必然大于父结点，我们可以按结点编号从大到小遍历（自底向上），对每个非根结点 $v$：
\[
G(v) = 1 + \bigoplus_{\substack{c \in \{0,\dots,25\} \\ \text{child}(v, c) \,\neq\, \text{null}}} G(\text{child}(v, c))
\]
叶结点无子结点，异或和为 $0$，故 $G(\text{leaf}) = 1$。

以样例为例验证公式。第一组样例：单词 \texttt{ab} 与 \texttt{cd}。Trie根有两个子结点 \texttt{a} 和 \texttt{c}，各为一条长度 $2$ 的链。叶结点 \texttt{b} 的 $G = 1$，\texttt{a} 的 $G = 1 + 1 = 2$；同理 \texttt{d} 的 $G = 1$，\texttt{c} 的 $G = 2$。总Grundy $= 2 \oplus 2 = 0$，后手必胜，输出 XiaoQiao，与样例一致。

第二组样例：单词 \texttt{cat} 与 \texttt{car}。叶结点 \texttt{t} 的 $G = 1$，叶结点 \texttt{r} 的 $G = 1$。结点 \texttt{a} 有两个子结点，$G(\texttt{a}) = 1 + (1 \oplus 1) = 1 + 0 = 1$。结点 \texttt{c} 有子结点 \texttt{a}（$G = 1$），$G(\texttt{c}) = 1 + 1 = 2$。根下仅一个子结点 \texttt{c}，总Grundy $= 2 \ne 0$，先手必胜，输出 XiaoLan，也与样例一致。从博弈角度看，小蓝第一步选择 \texttt{c} 即可清空整棵树，小桥无子可走。

\subsection{复杂度分析}

建Trie和自底向上计算Grundy值均仅对每个结点访问常数次。Trie的总结点数为所有单词的字符总数，即不超过 $10^6$。每个结点的Grundy计算只需遍历26个子结点并做异或，总操作量 $\mathrm{O}(|\Sigma| \cdot M)$，其中 $|\Sigma| = 26$，$M$ 为总结点数。整体时间复杂度 $\mathrm{O}(\sum |s_i|)$，空间复杂度同为 $\mathrm{O}(\sum |s_i|)$。对总字符数 $10^6$ 的输入，所有操作在毫秒级即可完成。

\subsection{实现细节}

在代码实现中，Trie使用扁平数组存储而非动态分配：\mil{nxt[node * 26 + c]} 表示从当前结点沿字符 $c$ 走到的子结点编号，\mil{-1} 表示无边。该设计避免了大量指针分配，访问速度更快，也使得结点按自然插入顺序获得递增编号，天然满足子结点编号大于父结点的性质。

自底向上遍历时，从最大编号结点递减至根结点（编号1）。根结点（编号0）对应空前缀，不参与Grundy值计算；其各子结点的Grundy值异或即为总Grundy值。总Grundy非零时先手小蓝必胜，否则后手小桥必胜。

\subsection{代码实现}
\inputminted{c++}{../src/p08_magic_prefix.cpp}
